Die Jacobimatrix von $\varphi$ ist allgemein

\mathdisp {\begin{pmatrix}  yz^2 &  xz^2 &  2xyz \\  y  &  x  &  -3z^2 \end{pmatrix}} { . }

Für den Punkt

\mavergleichskette
{\vergleichskette
{ P
}
{ = }{ (1,3,5)
}
{  }{
}
{    }{
}
{   }{
}
}
{}{}{}
liegt daher die Jacobimatrix

\mathdisp {\begin{pmatrix}  75 &  25 &  30 \\  3 &  1 &  -75  \end{pmatrix}} {  }

vor. Diese Matrix beschreibt die lineare Abbildung
\maabbdisp {L} {T_P\R^3 \cong \R^3} { T_{\varphi(P)} \R^2 \cong \R^2
} {}
bezüglich der Standardbasen. Wir bestimmen die Matrixdarstellung für die Abbildung
\maabbdisp {\bigwedge^2 L} {\bigwedge^2 \R^3 } { \bigwedge^2 \R^2
} {}
bezüglich der Basen
$e_1 \wedge e_2,\, e_1 \wedge e_3$ , $e_2 \wedge e_3$
\zusatzklammer {links} {} {}
und
\mathl{e_1 \wedge e_2}{}
\zusatzklammer {rechts} {} {.}
Dazu müssen wir die Bilder dieser Dachprodukte ausrechnen. Es ist

\mavergleichskettealign
{\vergleichskettealign
{ (\bigwedge^2 L ) (e_1 \wedge e_2)
}
{ =} { L(e_1) \wedge L(e_2)
}
{ =} { (75e_1 + 3e_2) \wedge (  25 e_1 + 1 e_2 )
}
{ =} { 75 e_1 \wedge e_2 + 75 e_2 \wedge  e_1
}
{ =} { 75 e_1 \wedge e_2 - 75 e_1 \wedge  e_2
}
}
{
\vergleichskettefortsetzungalign
{ =} { 0
}
{ } {}
{ } {}
{ } {}
}
{}{,}

\mavergleichskettealign
{\vergleichskettealign
{ (\bigwedge^2 L ) (e_1 \wedge e_3)
}
{ =} { L(e_1) \wedge L(e_3)
}
{ =} { (75e_1 + 3e_2) \wedge (  30 e_1 - 75 e_2 )
}
{ =} { - 75^2 e_1 \wedge e_2 + 90 e_2 \wedge  e_1
}
{ =} { (-5625 - 90) e_1 \wedge e_2
}
}
{
\vergleichskettefortsetzungalign
{ =} { -5715 e_1 \wedge e_2
}
{ } {}
{ } {}
{ } {}
}
{}{,}
und

\mavergleichskettealign
{\vergleichskettealign
{ (\bigwedge^2 L ) (e_2 \wedge e_3)
}
{ =} { L(e_2) \wedge L(e_3)
}
{ =} { (  25 e_1 + 1 e_2 )   \wedge ( 30 e_1 - 75 e_2)
}
{ =} { - 1875 e_1 \wedge e_2 + 30 e_2 \wedge  e_1
}
{ =} { -  1905 e_1 \wedge  e_2
}
}
{}
{}{.}
Die beschreibende Matrix ist also

\mathdisp {\left(  0,-5715,-1905                     \right)} { . }

 [[Kategorie:Latexseite]]
